\subsection{后继轨道、双横截面末端正规形与奇异环乘法语义的共同有限化障碍}\label{subsec:conclusion-successor-semiconjugacy-ledger-obstruction}
沿用定义 \ref{def:successor-fold}、定理 \ref{thm:successor-structure}、推论 \ref{cor:add-from-successor}、定理 \ref{thm:mul-from-successor}、定理 \ref{thm:pom-double-transversal-normal-form}、推论 \ref{cor:pom-projection-budget}、定理 \ref{thm:leyang-chebyshev-semiconjugacy}、推论 \ref{cor:leyang-endpoint-preimage-cos2-spectrum}、命题 \ref{prop:leyang-ellipse-layer-positive-real-crossing}、定理 \ref{thm:killo-prime-freedom-non-finitizability}、定理 \ref{thm:killo-godel-compression-not-finite-rank-homologizable} 与定理 \ref{thm:killo-localization-circle-dimension-prime-ledger-splitting} 的记号。前述结果已经分别给出：稳定地址层上的后继去原语化、包含商余在内的双横截面末端正规形、Lee--Yang 奇异环上的 Chebyshev 半共轭、以及有限生成加法账本对乘法语义的结构性阻断。把这些接口并读后，可以把“后继--折叠”与“幂映射--Chebyshev”压缩为同一乘法半群的两类忠实表示，并进一步得到它们在有限秩加法线性化障碍上的严格同型。

\begin{theorem}[稳定地址层的后继轨道正规形]\label{thm:conclusion-stable-successor-orbit-normal-form}
\leanverified{paper\_conclusion\_stable\_successor\_orbit\_normal\_form}
对任意 \(c\in X_\infty\)，有唯一表示
\[
c=S^{\circ\,\mathrm{Val}(c)}(0_X).
\]
因此，\((X_\infty,S)\) 与自然数后继系统 \((\NN,n\mapsto n+1)\) 在 \(\mathrm{Val}\) 下严格同构。
\end{theorem}
\begin{proof}
由定理 \ref{thm:successor-structure}，
\[
\mathrm{Val}(0_X)=0,\qquad
\mathrm{Val}(S(c))=\mathrm{Val}(c)+1,
\]
且 \(\mathrm{Val}\) 为双射。于是对任意 \(n\ge 0\) 归纳得到
\[
\mathrm{Val}\bigl(S^{\circ n}(0_X)\bigr)=n.
\]
取 \(n=\mathrm{Val}(c)\)，再由 \(\mathrm{Val}\) 的单射性可得
\[
c=S^{\circ\,\mathrm{Val}(c)}(0_X).
\]
唯一性同样来自 \(\mathrm{Val}\) 的单射性。证毕。
\end{proof}

\begin{corollary}[加法作为同一后继轨道上的平移读出]\label{cor:conclusion-stable-successor-translation-addition}
\leanverified{paper\_conclusion\_stable\_successor\_translation\_addition}
对任意 \(c,d\in X_\infty\)，有
\[
c\oplus d=S^{\circ\,\mathrm{Val}(d)}(c)=\Fold_\infty(c+d).
\]
因此，稳定加法只是后继轨道上的位移读出，而不是独立原语。
\end{corollary}
\begin{proof}
由推论 \ref{cor:add-from-successor}，
\[
c\oplus d=S^{\circ\,\mathrm{Val}(d)}(c).
\]
再由推论 \ref{cor:add-as-fold}，
\[
c\oplus d=\Fold_\infty(c+d).
\]
证毕。
\end{proof}

\begin{theorem}[算术接口的双横截面终端正规形]\label{thm:conclusion-arithmetic-double-transversal-terminal-normal-form}
\leanverified{paper\_conclusion\_arithmetic\_double\_transversal\_terminal\_normal\_form}
对任意配置对象 \(a\in\mathcal D\) 与任意整数 \(b>0\)，存在唯一三元组
\[
(x,q,r)\in X_\infty\times\ZZ\times\mathcal F_b
\]
使得
\[
x=\Fold_\infty(a),\qquad
r=\Pi_b(\mathrm{Val}(a)),\qquad
q=\frac{\mathrm{Val}(a)-r}{b},
\]
并满足
\[
\mathrm{Val}(a)=\mathrm{Val}(x)=qb+r.
\]
\end{theorem}
\begin{proof}
这正是定理 \ref{thm:pom-double-transversal-normal-form} 的内容。证毕。
\end{proof}

\begin{corollary}[末端投影的两轴完备性]\label{cor:conclusion-arithmetic-terminal-two-axis-completeness}
\leanverified{paper\_conclusion\_arithmetic\_terminal\_two\_axis\_completeness}
任意由 \(\oplus,\otimes,\mathrm{div},\mathrm{quo},\mathrm{rem}\) 复合得到的表达式，都存在等价实现，使整条求值链中
\[
\Fold_\infty \text{ 至多调用一次},\qquad
\Pi_b \text{ 至多调用一次},
\]
且二者都可推迟到末端。特别地，一旦给定 \((x,r)\)，则
\[
q=\frac{\mathrm{Val}(x)-r}{b},
\]
故真正不可消去的外部读出仅剩值纤维横截面与同余横截面两条轴。
\end{corollary}
\begin{proof}
推论 \ref{cor:pom-projection-budget} 给出两类投影都可推迟到末端且至多调用一次。另一方面，由定理 \ref{thm:conclusion-arithmetic-double-transversal-terminal-normal-form}，
\[
\mathrm{Val}(x)=qb+r,
\]
故
\[
q=\frac{\mathrm{Val}(x)-r}{b}.
\]
因此 \((q,r)\) 中的商数部分并不构成独立寄存器。证毕。
\end{proof}

\begin{theorem}[稳定乘法作用的忠实乘法半群表示]\label{thm:conclusion-stable-successor-faithful-semigroup}
记 \(Z:=\mathrm{Val}^{-1}\)，并对每个 \(n\ge 1\) 定义
\[
M_n:X_\infty\to X_\infty,\qquad
M_n(c):=Z\!\bigl(n\,\mathrm{Val}(c)\bigr).
\]
则对任意 \(c,d\in X_\infty\) 有
\[
c\otimes d=M_{\mathrm{Val}(d)}(c)=S^{\circ\,\mathrm{Val}(c)\mathrm{Val}(d)}(0_X).
\]
此外，映射
\[
\NN_{\ge 1}\longrightarrow \operatorname{End}(X_\infty),\qquad
n\longmapsto M_n
\]
在复合下给出乘法半群 \((\NN_{\ge 1},\cdot)\) 的忠实表示：
\[
M_m\circ M_n=M_{mn},
\qquad
M_m=M_n\Longrightarrow m=n.
\]
\leanverified{paper\_conclusion\_stable\_successor\_faithful\_semigroup}
\end{theorem}
\begin{proof}
由稳定乘法的定义，
\[
c\otimes d
=
Z\!\bigl(\mathrm{Val}(c)\mathrm{Val}(d)\bigr)
=
M_{\mathrm{Val}(d)}(c).
\]
再由定理 \ref{thm:conclusion-stable-successor-orbit-normal-form}，
\[
Z(k)=S^{\circ k}(0_X)\qquad(k\in\NN),
\]
故
\[
c\otimes d
=
S^{\circ\,\mathrm{Val}(c)\mathrm{Val}(d)}(0_X).
\]
对半群关系，只需计算
\[
\mathrm{Val}\bigl((M_m\circ M_n)(c)\bigr)
=
m\,\mathrm{Val}(M_n(c))
=
mn\,\mathrm{Val}(c)
=
\mathrm{Val}(M_{mn}(c)).
\]
由 \(\mathrm{Val}\) 的单射性得 \(M_m\circ M_n=M_{mn}\)。
若 \(M_m=M_n\)，取 \(1_X:=S(0_X)\)，则
\[
m=\mathrm{Val}(M_m(1_X))=\mathrm{Val}(M_n(1_X))=n,
\]
故表示忠实。证毕。
\end{proof}

\begin{theorem}[奇异环 Chebyshev 半共轭的忠实乘法半群表示]\label{thm:conclusion-singular-ring-faithful-semigroup}
\leanverified{paper\_conclusion\_singular\_ring\_faithful\_semigroup}
在定理 \ref{thm:leyang-chebyshev-semiconjugacy} 的记号下，映射
\[
\NN_{\ge 1}\longrightarrow \QQ(t),\qquad
n\longmapsto f_n
\]
在函数复合下给出 \((\NN_{\ge 1},\cdot)\) 的忠实表示：
\[
f_m\circ f_n=f_{mn},
\qquad
f_m=f_n\Longrightarrow m=n.
\]
\end{theorem}
\begin{proof}
半群关系 \(f_m\circ f_n=f_{mn}\) 已由定理 \ref{thm:leyang-chebyshev-semiconjugacy} 给出。
由该定理中的显式公式
\[
f_n(t)=-\frac{1}{4\,P_n\!\left(-\frac{1}{4t}\right)},
\qquad
\deg P_n=n,
\]
若记
\[
\widetilde P_n(t):=t^nP_n\!\left(-\frac{1}{4t}\right)\in \QQ[t],
\]
则 \(\widetilde P_n\) 的次数为 \(n\)，且有
\[
f_n(t)=-\frac{t^n}{4\,\widetilde P_n(t)}.
\]
因此 \(f_n\) 作为 \(\PP^1\) 上的有理自映射，其次数恰为 \(n\)。若 \(f_m=f_n\)，则其次数相等，遂 \(m=n\)。证毕。
\end{proof}

\begin{corollary}[复合不可约元与算术素数的一致性]\label{cor:conclusion-singular-ring-prime-composition-irreducibility}
\leanverified{paper\_conclusion\_singular\_ring\_prime\_composition\_irreducibility}
对任意 \(n\ge 2\)，有理映射 \(f_n\) 在函数复合半群中不可约，当且仅当整数 \(n\) 为素数。
\end{corollary}
\begin{proof}
若 \(n=ab\) 且 \(a,b>1\)，则由定理 \ref{thm:conclusion-singular-ring-faithful-semigroup}，
\[
f_n=f_{ab}=f_a\circ f_b,
\]
故 \(f_n\) 可约。
反之，若 \(f_n=f_u\circ f_v\) 且 \(u,v>1\)，则同一定理给出
\[
f_n=f_{uv}.
\]
由忠实性立得 \(n=uv\)，故 \(n\) 非素。证毕。
\end{proof}

\begin{theorem}[奇异端点原像塔的紧化乘法律]\label{thm:conclusion-singular-ring-endpoint-preimage-tower}
\leanverified{paper\_conclusion\_singular\_ring\_endpoint\_preimage\_tower}
对任意 \(n\ge 2\)，有
\[
f_n^{-1}\!\left(-\frac14\right)
=
\left\{
-\frac{1}{4\cos^2(\pi k/n)}:\ 1\le k\le n-1
\right\},
\]
其中当 \(n\) 为偶数且 \(k=n/2\) 时，对应紧化端点 \(-\infty\)。该原像集的不同点个数恰为
\[
\left\lfloor\frac{n}{2}\right\rfloor.
\]
并且对任意 \(m,n\ge 1\)，成立集合恒等式
\[
f_m^{-1}\!\left(f_n^{-1}\!\left(-\frac14\right)\right)
=
f_{mn}^{-1}\!\left(-\frac14\right).
\]
\end{theorem}
\begin{proof}
原像描述由推论 \ref{cor:leyang-endpoint-preimage-cos2-spectrum} 直接给出。
对计数部分，注意
\[
\cos^2\!\frac{\pi k}{n}=\cos^2\!\frac{\pi(n-k)}{n},
\]
故 \(k\) 与 \(n-k\) 配对给出同一点。若 \(n\) 为奇数，得到 \((n-1)/2\) 个不同有限点；若 \(n\) 为偶数，则除配对点外还多出 \(k=n/2\) 所对应的紧化端点 \(-\infty\)，总数为 \(n/2\)。统一写成 \(\lfloor n/2\rfloor\)。

再由定理 \ref{thm:conclusion-singular-ring-faithful-semigroup}，
\[
f_m^{-1}\!\left(f_n^{-1}\!\left(-\frac14\right)\right)
=
(f_n\circ f_m)^{-1}\!\left(-\frac14\right)
=
f_{nm}^{-1}\!\left(-\frac14\right)
=
f_{mn}^{-1}\!\left(-\frac14\right).
\]
证毕。
\end{proof}

\begin{theorem}[稳定乘法语义的有限秩账本不可能性]\label{thm:conclusion-stable-multiplication-finite-rank-ledger-impossibility}
\leanverified{paper\_conclusion\_stable\_multiplication\_finite\_rank\_ledger\_impossibility}
记
\[
X_\infty^{\times}:=\{c\in X_\infty:\ \mathrm{Val}(c)\ge 1\}.
\]
设 \(e:X_\infty^{\times}\to\ZZ\) 为任意单射。则不存在有限生成交换群 \(L\) 与映射
\[
s:X_\infty^{\times}\to L
\]
使得
\[
\Phi:X_\infty^{\times}\to \ZZ\oplus L,\qquad
\Phi(c):=(e(c),s(c))
\]
既为单射，又满足
\[
\Phi(c\otimes d)=\Phi(c)+\Phi(d)
\qquad
(\forall\,c,d\in X_\infty^{\times}).
\]
\end{theorem}
\begin{proof}
反设存在这样的 \((L,s)\)。由定理 \ref{thm:mul-definitional}，
\[
\mathrm{Val}:(X_\infty,\oplus,\otimes)\longrightarrow (\NN,+,\times)
\]
是半环同构；特别地，限制后得到双射
\[
Z:=\mathrm{Val}^{-1}:\NN_{\ge 1}\longrightarrow X_\infty^{\times}.
\]
于是可定义
\[
\widetilde\Phi:\NN_{\ge 1}\to \ZZ\oplus L,
\qquad
\widetilde\Phi(n):=\Phi(Z(n)).
\]
由于 \(\Phi\) 与 \(Z\) 都单射，\(\widetilde\Phi\) 亦单射。再对任意 \(m,n\ge 1\)，有
\[
\widetilde\Phi(mn)
=
\Phi\bigl(Z(mn)\bigr)
=
\Phi\bigl(Z(m)\otimes Z(n)\bigr)
=
\Phi\bigl(Z(m)\bigr)+\Phi\bigl(Z(n)\bigr)
=
\widetilde\Phi(m)+\widetilde\Phi(n).
\]
故 \(\widetilde\Phi\) 给出 \(\NN_{\times}\) 到 \((\ZZ\oplus L,+)\) 的单射幺半群同态。
这与定理 \ref{thm:killo-godel-compression-not-finite-rank-homologizable} 矛盾。证毕。
\end{proof}

\begin{theorem}[单位模壳是一维奇异射线束缚的极大壳]\label{thm:conclusion-singular-ring-unit-circle-maximal-shell}
\leanverified{paper\_conclusion\_singular\_ring\_unit\_circle\_maximal\_shell}
令
\[
\tau(z):=J(z^2)=-\frac{1}{(z+z^{-1})^2}.
\]
则当 \(|z|=1\) 时，总有
\[
\tau(z)\in\left(-\infty,-\frac14\right].
\]
若固定 \(|z|=r\neq 1\) 并令 \(z=re^{\mathrm{i}\theta}\) 变化，则 \(\tau(z)\) 的像同时与负实轴与正实轴相交。特别地，单位模约束恰是逆平方门保持一维奇异射线退化的最大壳层；任意径向脱离都会立即打破这一一维束缚。
\end{theorem}
\begin{proof}
当 \(|z|=1\) 时，结论由命题 \ref{prop:jouwkowsky-inverse-square} 的第二条给出。
当 \(|z|=r\neq 1\) 时，命题 \ref{prop:leyang-ellipse-layer-positive-real-crossing} 说明
\[
\theta=0:\quad \tau(z)=-\frac{1}{(r+r^{-1})^2}<0,
\qquad
\theta=\frac{\pi}{2}:\quad \tau(z)=\frac{1}{(r-r^{-1})^2}>0.
\]
故其像同时穿过负实轴与正实轴。证毕。
\end{proof}

\begin{theorem}[稳定乘法作用与奇异环半共轭的共同有限加法线性化障碍]\label{thm:conclusion-common-finite-additive-linearization-obstruction}
\leanverified{paper\_conclusion\_common\_finite\_additive\_linearization\_obstruction}
不存在任何有限 \(k\) 与单射幺半群同态，把下列任一忠实乘法表示线性化到 \((\NN^k,+)\)：
\[
\{M_n\}_{n\ge 1}\subset \operatorname{End}(X_\infty),
\qquad
\{f_n\}_{n\ge 1}\subset \QQ(t).
\]
更一般地，不存在任何有限生成交换可消去幺半群容纳这两类表示中的任一者作为忠实加法像；同样也不存在有限生成交换群账本 \(L\) 使其可忠实同态化到 \((\ZZ\oplus L,+)\)。
\end{theorem}
\begin{proof}
由定理 \ref{thm:conclusion-stable-successor-faithful-semigroup}，\(\{M_n\}_{n\ge 1}\) 与 \((\NN_{\ge 1},\cdot)\) 忠实同构；由定理 \ref{thm:conclusion-singular-ring-faithful-semigroup}，\(\{f_n\}_{n\ge 1}\) 也与 \((\NN_{\ge 1},\cdot)\) 忠实同构。
若其中任一者可单射同态进 \((\NN^k,+)\) 或更一般的有限生成交换可消去幺半群 \(M\)，则与上述忠实同构复合后便得到
\[
\NN_{\times}\hookrightarrow M,
\]
这与定理 \ref{thm:killo-prime-freedom-non-finitizability} 矛盾。
若进一步存在有限生成交换群 \(L\) 使其可忠实同态化到 \((\ZZ\oplus L,+)\)，则同样得到
\[
\NN_{\times}\hookrightarrow (\ZZ\oplus L,+),
\]
这被定理 \ref{thm:killo-godel-compression-not-finite-rank-homologizable} 所排除。证毕。
\end{proof}

\begin{corollary}[有限秩连续补全的单圆可见轴与 prime-ledger 核]\label{cor:conclusion-finite-rank-completion-single-circle-prime-ledger}
设 \(S\subset\PP\) 为有限素数集，并记 \(G_S=\ZZ[S^{-1}]\)。则其 Pontryagin 对偶满足短正合列
\[
0\longrightarrow \prod_{p\in S}\ZZ_p
\longrightarrow \widehat{G_S}
\longrightarrow \TT
\longrightarrow 0,
\qquad
\cdim_{\mathrm{fr}}(G_S)=1.
\]
因此，一切有限秩连续补全至多留下一个可见圆维 \(\TT\)；全部真正的乘法/局部化残差都被迫滞留在全不连通的 prime-ledger 核
\[
\prod_{p\in S}\ZZ_p
\]
中，并且 \(S\) 可由该核的 \(p\)-主因子完全恢复。
\end{corollary}
\begin{proof}
这正是定理 \ref{thm:killo-localization-circle-dimension-prime-ledger-splitting} 的内容。证毕。
\end{proof}

\begin{conjecture}[稳定乘法作用到奇异射线动力学的谱函子]\label{conj:conclusion-stable-successor-singular-ring-spectral-functor}
预期存在一个由投影词正规形、primitive 读出与末端双横截面共同组织的谱函子
\[
\mathfrak F
\]
把稳定地址层上的乘法作用
\[
M_n:X_\infty\to X_\infty
\]
送到奇异环上的 Chebyshev 半共轭作用
\[
f_n:\tau(\TT)\to \tau(\TT),
\]
并满足
\[
\mathfrak F(M_m\circ M_n)=\mathfrak F(M_m)\circ \mathfrak F(M_n),
\qquad
\mathfrak F(M_n)=f_n.
\]
若该函子得以显式构造，则稳定算术中的“后继--折叠”与奇异环动力学中的“幂映射--Chebyshev”将不再只是两条并行语义，而会被识别为同一乘法半群在对象层与谱层上的两种正交投影。
\end{conjecture}
